%───────────────────────────────────────────────────────────────────────────────
\chapter{Estructura del documento}
\label{ch:estructura}
%───────────────────────────────────────────────────────────────────────────────

Un documento académico de largo aliento tiene estructura en múltiples escalas:
la jerarquía visual de títulos que orienta al lector, el sistema de referencias
cruzadas que da cohesión al argumento, el frontmatter que contextualiza la obra,
y la bibliografía que ancla cada afirmación. Este capítulo documenta cómo el
template gestiona cada una de esas escalas y cómo usarlas correctamente.

%───────────────────────────────────────────────────────────────────────────────
\section{Jerarquía de títulos}
\label{sec:titulos}
%───────────────────────────────────────────────────────────────────────────────

El template configura cuatro niveles de título mediante
\texttt{titlesec}~\cite{titlesec2021}, cada uno con un color semántico
distinto definido en la paleta bajo el prefijo \texttt{doc-*}:

\begin{acadtable}[H]{tab:titulo-niveles}{Jerarquía de títulos y sus colores}
  \begin{tabularx}{\textwidth}{l l l X}
    \toprule
    \headerrow
    \thead{Comando} & \thead{Alias de color} & \thead{Color Ctp} & \thead{Uso} \\
    \midrule
    \texttt{\textbackslash chapter}
      & \texttt{doc-chapter} & Pink
      & Unidad temática mayor \\
    \texttt{\textbackslash section}
      & \texttt{doc-section} & Sky
      & Subtema dentro del capítulo \\
    \texttt{\textbackslash subsection}
      & \texttt{doc-subsection} & Lavender
      & División de una sección \\
    \texttt{\textbackslash subsubsection}
      & \texttt{doc-subsubsection} & Peach
      & Detalle dentro de subsección \\
    \bottomrule
  \end{tabularx}
\end{acadtable}

La progresión de colores no es arbitraria: contraste decreciente con la
profundidad. El capítulo ancla visualmente la página con el Pink más
prominente; la subsubsección aporta apenas un tono Peach que diferencia
sin competir. El lector percibe la jerarquía antes de leer.

\begin{observacion}{Cambiar colores de títulos}
Los colores se modifican exclusivamente en
\texttt{themes/catppuccin-palette.tex}, cambiando el alias correspondiente.
Por ejemplo, para que las secciones usen Mauve en lugar de Sky:
\texttt{\textbackslash colorlet\{doc-section\}\{CtpMauve\}}.
Nunca editar \texttt{config/titles.tex} para poner colores directamente.
\end{observacion}

\subsection{Títulos sin numeración}
\label{subsec:titulos-nonumber}

Los comandos con asterisco (\texttt{\textbackslash section*},
\texttt{\textbackslash subsection*}) suprimen la numeración pero
mantienen el color y la tipografía del nivel correspondiente. Son útiles
para secciones estructurales como prefacios o glosarios que no deben
aparecer numeradas en el texto pero sí en la tabla de contenidos:

\begin{lstlisting}[style=latex, frame=none, numbers=none]
\section*{Notación}
\addcontentsline{toc}{section}{Notación}  % añadir al TOC manualmente
\end{lstlisting}

\subsubsection{Párrafos como cuarto nivel}
\label{subsubsec:parrafos}

El cuarto nivel visible en el texto es \texttt{\textbackslash subsubsection}.
Existe un quinto nivel, \texttt{\textbackslash paragraph}, que aparece como
texto inline sin salto de línea propio. Su uso debe ser esporádico: si
necesitas un quinto nivel estructural con frecuencia, la estructura del
documento probablemente necesita reorganización.

%───────────────────────────────────────────────────────────────────────────────
\section{Sistema de referencias cruzadas}
\label{sec:referencias-cruzadas}
%───────────────────────────────────────────────────────────────────────────────

El template usa \texttt{cleveref}~\cite{cleveref2022} para referencias
inteligentes. A diferencia de \texttt{\textbackslash ref}, que produce
solo el número, \texttt{\textbackslash cref} produce el tipo y el número
correctos en el idioma configurado:

\begin{definicion}{Referencias cruzadas con \texttt{cleveref}}
\texttt{\textbackslash cref\{etiqueta\}} produce la referencia completa
en minúscula: \textit{ec.~(2.1)}, \textit{fig.~3.2}, \textit{def.~1.4}.
\texttt{\textbackslash Cref\{etiqueta\}} produce la versión capitalizada
para inicio de oración: \textit{Ecuación~(2.1)}, \textit{Figura~3.2}.
Para rangos, \texttt{\textbackslash crefrange\{etiq1\}\{etiq2\}} produce
\textit{ecs.~(2.1)--(2.3)} automáticamente.
\end{definicion}

\subsection{Convención de etiquetas}
\label{subsec:etiquetas}

Una convención consistente de prefijos hace el código fuente legible y
evita colisiones entre etiquetas de distintos archivos:

\begin{acadtable}[H]{tab:label-prefijos}{Prefijos de etiquetas recomendados}
  \begin{tabularx}{\textwidth}{l l X}
    \toprule
    \headerrow
    \thead{Tipo} & \thead{Prefijo} & \thead{Ejemplo} \\
    \midrule
    Capítulo      & \texttt{ch:}   & \texttt{ch:estructura} \\
    Sección       & \texttt{sec:}  & \texttt{sec:referencias-cruzadas} \\
    Subsección    & \texttt{subsec:} & \texttt{subsec:etiquetas} \\
    Ecuación      & \texttt{eq:}   & \texttt{eq:schrodinger} \\
    Figura        & \texttt{fig:}  & \texttt{fig:convergencia} \\
    Tabla         & \texttt{tab:}  & \texttt{tab:label-prefijos} \\
    Algoritmo     & \texttt{alg:}  & \texttt{alg:gradient} \\
    Código        & \texttt{lst:}  & \texttt{lst:perceptron} \\
    Apéndice      & \texttt{ap:}   & \texttt{ap:galeria} \\
    Definición    & \texttt{def:}  & \texttt{def:hilbert} \\
    Teorema       & \texttt{thm:}  & \texttt{thm:riesz} \\
    \bottomrule
  \end{tabularx}
\end{acadtable}

\subsection{Referencias a entornos propios}
\label{subsec:cref-entornos}

Los entornos teoremáticos del template crean contadores propios. Para que
\texttt{cleveref} los reconozca, sus contadores están registrados en
\texttt{config/packages.tex}. El patrón de etiqueta es consistente con
los prefijos de la \cref{tab:label-prefijos}: poner el
\texttt{\textbackslash label} dentro del entorno, después del título:

\begin{lstlisting}[style=latex, frame=none, numbers=none]
\begin{definicion}{Espacio de Hilbert}
\label{def:hilbert}
Un espacio de Hilbert es...
\end{definicion}

% Más adelante en el texto:
Como se establece en la \cref{def:hilbert}, un espacio de Hilbert...
\end{lstlisting}

\begin{observacion}{Posición del \texttt{\textbackslash label}}
En entornos de \texttt{tcolorbox}, el \texttt{\textbackslash label} debe
ir \textit{dentro} del entorno, no en el argumento del título. El título
es un argumento de texto, no un entorno flotante. Poner el label fuera
produce un anchor de hyperref incorrecto.
\end{observacion}

%───────────────────────────────────────────────────────────────────────────────
\section{Frontmatter}
\label{sec:frontmatter}
%───────────────────────────────────────────────────────────────────────────────

El frontmatter agrupa las páginas previas al contenido. El template lo
divide en tres archivos en \texttt{frontmatter/}, cada uno con una
responsabilidad única:

\paragraph{\texttt{titlepage.tex} --- Portada.}
Genera la portada completa a partir de las variables de \texttt{metadata.tex}.
No requiere edición directa salvo cambios de diseño estructural. Los contenidos
se controlan exclusivamente desde \texttt{metadata.tex}.

\paragraph{\texttt{preliminary.tex} --- Páginas preliminares.}
Muestra dedicatoria, epígrafe y agradecimientos en páginas separadas.
Cada sección es condicional: si la variable correspondiente está vacía
(\texttt{\textbackslash docDedicatoria\{\}}), la página no se genera.
Esto permite activar o desactivar secciones sin editar el archivo.

\paragraph{\texttt{preface.tex} --- Prefacio.}
Capítulo sin numerar (\texttt{\textbackslash chapter*}) añadido al TOC
manualmente. A diferencia de los preliminares, el prefacio tiene contenido
narrativo propio --- motivación del trabajo, agradecimientos extendidos,
contexto histórico --- y se edita directamente como cualquier capítulo.

\begin{observacion}{Paginación del frontmatter}
Por convención tipográfica, el frontmatter usa numeración romana
(\texttt{\textbackslash pagenumbering\{roman\}}) y el cuerpo del documento
numeración arábiga (\texttt{\textbackslash pagenumbering\{arabic\}}).
El template aplica esto en \texttt{main.tex}. Si quieres que el TOC empiece
en la página~i, asegúrate de que el cambio de numeración ocurra antes de
\texttt{\textbackslash tableofcontents}.
\end{observacion}

%───────────────────────────────────────────────────────────────────────────────
\section{Tabla de contenidos}
\label{sec:toc}
%───────────────────────────────────────────────────────────────────────────────

\texttt{\textbackslash tableofcontents} genera el TOC automáticamente a
partir de los títulos del documento. El nivel de profundidad se controla
con \texttt{tocdepth}:

\begin{lstlisting}[style=latex, frame=none, numbers=none]
\setcounter{tocdepth}{2}    % muestra hasta \subsection
\setcounter{tocdepth}{3}    % muestra hasta \subsubsection (default)
\end{lstlisting}

El template configura \texttt{secnumdepth} en 3 por defecto --- todos los
niveles hasta \texttt{\textbackslash subsubsection} están numerados. Si una
sección con asterisco debe aparecer en el TOC (como el prefacio), se añade
manualmente:

\begin{lstlisting}[style=latex, frame=none, numbers=none]
\chapter*{Prefacio}
\addcontentsline{toc}{chapter}{Prefacio}
\end{lstlisting}

%───────────────────────────────────────────────────────────────────────────────
\section{Bibliografía por capítulo}
\label{sec:bibliografia}
%───────────────────────────────────────────────────────────────────────────────

\begin{definicion}{Refsection de \texttt{biblatex}}
Con la opción \texttt{refsection=chapter}, \texttt{biblatex} trata cada
\texttt{\textbackslash chapter} (o cada \texttt{\textbackslash include})
como una sección bibliográfica independiente. Las citas de un capítulo
solo aparecen en la bibliografía de ese capítulo.
\texttt{\textbackslash printbibliography[heading=subbibliography]} al
final de cada capítulo imprime solo las referencias citadas en ese capítulo,
con numeración reiniciada.
\end{definicion}

El patrón a seguir en cada capítulo nuevo es:

\begin{lstlisting}[style=latex, frame=none, numbers=none]
% 1. Crear references/chapterN.bib con las referencias del capitulo
% 2. En metadata.tex, anadir:
\addbibresource{references/chapterN.bib}

% 3. Al final del archivo chapterN.tex:
\printbibliography[heading=subbibliography]
\end{lstlisting}

\begin{observacion}{Referencias compartidas entre capítulos}
Si una referencia aparece en varios capítulos, debe estar en el
\texttt{.bib} del capítulo donde se cita por primera vez, o en el
\texttt{references.bib} global. \texttt{biblatex} admite múltiples
\texttt{\textbackslash addbibresource} y no duplica entradas aunque
la misma clave aparezca en varios archivos.
\end{observacion}

%───────────────────────────────────────────────────────────────────────────────
\section{Apéndices}
\label{sec:apendices}
%───────────────────────────────────────────────────────────────────────────────

Los apéndices~\cite{appendix2020} usan el entorno \texttt{appendices}
provisto por el paquete homónimo. Dentro de él, cada apéndice es un
\texttt{\textbackslash chapter} normal cuya numeración cambia
automáticamente a letras (A, B, C\ldots):

\begin{lstlisting}[style=latex, frame=none, numbers=none]
% En main.tex, después de los capítulos:
\begin{appendices}
  %───────────────────────────────────────────────────────────────────────────────
\chapter{Style Guide de Catppuccin Académico}
\label{ap:styleguide}
%───────────────────────────────────────────────────────────────────────────────

Este apéndice es el documento normativo del sistema de color del template.
Define la arquitectura, las reglas de decisión, el inventario completo de
aliases y el protocolo de extensión. Todo color nuevo que se añada al
template debe pasar por este documento primero.

%───────────────────────────────────────────────────────────────────────────────
\section{Arquitectura de tres capas}
\label{ap:sec:arquitectura}
%───────────────────────────────────────────────────────────────────────────────

\begin{definicion}{Las tres capas del sistema de color}
\label{ap:def:tres-capas}
El sistema de color opera en tres niveles de abstracción con una regla
de visibilidad estricta:

\textbf{Nivel 0 — Paleta cruda (\texttt{Ctp*}):} Valores hexadecimales
absolutos del flavor activo. Existen 26 colores: \texttt{CtpRosewater},
\texttt{CtpFlamingo}, \texttt{CtpPink}, \texttt{CtpMauve}, \texttt{CtpRed},
\texttt{CtpMaroon}, \texttt{CtpPeach}, \texttt{CtpYellow}, \texttt{CtpGreen},
\texttt{CtpTeal}, \texttt{CtpSky}, \texttt{CtpSapphire}, \texttt{CtpBlue},
\texttt{CtpLavender}, \texttt{CtpText}, \texttt{CtpSubtext1},
\texttt{CtpSubtext0}, \texttt{CtpOverlay2}, \texttt{CtpOverlay1},
\texttt{CtpOverlay0}, \texttt{CtpSurface2}, \texttt{CtpSurface1},
\texttt{CtpSurface0}, \texttt{CtpBase}, \texttt{CtpMantle},
\texttt{CtpCrust}.

\textbf{Nivel 1 — Aliases estructurales:} Roles de UI genéricos sin
semántica de dominio. Ejemplos: \texttt{background}, \texttt{surface0},
\texttt{overlay0}, \texttt{maintext}. Describen propiedades visuales de
elementos de interfaz.

\textbf{Nivel 2 — Aliases semánticos:} Roles de dominio del documento.
Ejemplos: \texttt{env-def}, \texttt{doc-chapter}, \texttt{code-keyword}.
Describen la intención del elemento en el documento.
\end{definicion}

La regla de visibilidad es \textbf{absoluta}:

\begin{estrella}{Invariante de capas}
\label{ap:star:invariante}
\texttt{Ctp*} solo aparece en \texttt{themes/catppuccin-palette.tex}.\\
Los archivos de \texttt{config/} consumen aliases semánticos.\\
Los archivos de contenido (\texttt{chapters/}, \texttt{environments/})
consumen comandos del template.

Cuando este invariante se mantiene, cambiar el flavor del documento
entero es descomentar una línea en \texttt{catppuccin-palette.tex}.
Cuando se viola, el cambio de flavor requiere buscar y editar en
múltiples archivos.
\end{estrella}

%───────────────────────────────────────────────────────────────────────────────
\section{Árbol de decisión para asignar colores}
\label{ap:sec:decision}
%───────────────────────────────────────────────────────────────────────────────

Cuando se necesita un color nuevo en cualquier archivo de config o de
contenido, la secuencia de preguntas es:

\begin{enumerate}[label=\textbf{\arabic*.}, leftmargin=2em]
  \item \textbf{¿Estoy en \texttt{catppuccin-palette.tex}?}
    Si sí, puedo usar \texttt{Ctp*} libremente para definir aliases.
    Si no, nunca uso \texttt{Ctp*}: avanzo al paso 2.

  \item \textbf{¿Ya existe un alias semántico para este rol?}
    Buscar en la paleta con el prefijo del dominio (\texttt{env-},
    \texttt{code-}, \texttt{tikz-}, etc.). Si existe, lo consumo.
    Si no existe, avanzo al paso 3.

  \item \textbf{¿El color tiene semántica de dominio o es infraestructura?}
    Si es semántica de dominio (``el borde de un estado aceptador de
    autómata'', ``el color de una keyword''), añado un alias semántico
    nuevo en la paleta bajo el prefijo correspondiente.
    Si es infraestructura genérica (``un fondo de superficie'', ``texto
    secundario''), uso un alias estructural existente (\texttt{surface0},
    \texttt{subtext0}).

  \item \textbf{¿Es un cambio independiente de otros elementos?}
    Si el color de este elemento siempre cambia junto con otro, comparten
    alias. Si pueden cambiar independientemente, tienen aliases separados.
\end{enumerate}

\begin{observacion}{La regla de granularidad mínima}
Un elemento nuevo merece su propio alias semántico si y solo si su
cambio de color es \textit{independiente} de otros elementos. Si el
label de figura y el label de tabla siempre deben ser del mismo color,
comparten \texttt{caption-label}. Si pueden diferir (como es el caso en
este template), tienen \texttt{caption-fig} y \texttt{caption-tab}
separados.
\end{observacion}

%───────────────────────────────────────────────────────────────────────────────
\section{Convención de nombres}
\label{ap:sec:nombres}
%───────────────────────────────────────────────────────────────────────────────

\begin{definicion}{Convención kebab-case}
\label{ap:def:kebab}
Todos los aliases siguen la convención:
\[
  \texttt{\{dominio\}-\{elemento\}[-\{variante\}]}
\]
en kebab-case. El dominio identifica la categoría; el elemento describe
el rol dentro del dominio; la variante (\texttt{frame}, \texttt{back},
\texttt{text}) solo aparece cuando el elemento tiene múltiples roles de
color. Los aliases estructurales (Nivel 1) no llevan prefijo de dominio.
\end{definicion}

\begin{acadtable}[H]{ap:tab:dominios}{Dominios activos y sus prefijos}
  \begin{tabularx}{\textwidth}{l l X}
    \toprule
    \headerrow
    \thead{Prefijo} & \thead{Dominio} & \thead{Cubre} \\
    \midrule
    (sin prefijo)    & Infraestructura     & Fondos, superficies, texto base \\
    \texttt{doc-}    & Documento           & Jerarquía de títulos \\
    \texttt{link-}   & Hipervínculos       & Referencias, URLs, citas \\
    \texttt{env-}    & Entornos            & Definición, teorema, ejemplo, etc. \\
    \texttt{eq-}     & Ecuaciones          & Coloreado semántico de matemáticas \\
    \texttt{code-}   & Código fuente       & Listings, syntax highlighting \\
    \texttt{alg-}    & Algoritmos          & algorithm2e, pseudocódigo \\
    \texttt{plot-}   & Gráficas            & pgfplots, series, ejes \\
    \texttt{tikz-}   & Diagramas TikZ      & Nodos, aristas, autómatas \\
    \texttt{table-}  & Tablas              & Encabezados, filas, reglas \\
    \texttt{caption-}& Captions            & Labels por tipo de flotante \\
    \texttt{fb-}     & Feedback            & Éxito, error, advertencia, info \\
    \texttt{diff-}   & Diffs               & Cambios insertados/eliminados \\
    \texttt{rainbow-}& Secuencia           & N colores distinguibles ordenados \\
    \texttt{frontmatter-} & Portada        & Elementos de la página de título \\
    \bottomrule
  \end{tabularx}
\end{acadtable}

%───────────────────────────────────────────────────────────────────────────────
\section{Inventario completo de aliases}
\label{ap:sec:inventario}
%───────────────────────────────────────────────────────────────────────────────

\subsection{Infraestructura visual (Nivel 1)}
\label{ap:subsec:infraestructura}

\begin{acadtable}[H]{ap:tab:infra}{Aliases estructurales: fondos, superficies y tipografía}
  \begin{tabularx}{\textwidth}{l l X}
    \toprule
    \headerrow
    \thead{Alias} & \thead{Color Ctp} & \thead{Fundamento Style Guide} \\
    \midrule
    \texttt{background}   & Base      & Background Pane \\
    \texttt{secondarybg}  & Mantle    & Secondary Panes \\
    \texttt{crustbg}      & Crust     & Secondary Panes (más oscuro) \\
    \texttt{surface0/1/2} & Surface 0/1/2 & Surface Elements \\
    \texttt{overlay0/1/2} & Overlay 0/1/2 & Overlays \\
    \texttt{maintext}     & Text      & Body Copy \\
    \texttt{subtext0/1}   & Subtext 0/1 & Sub-Headlines, Labels \\
    \texttt{subtletext}   & Overlay 1 & Subtle — texto terciario \\
    \bottomrule
  \end{tabularx}
\end{acadtable}

\subsection{Estructura del documento (\texttt{doc-})}

\begin{acadtable}[H]{ap:tab:doc}{Aliases de jerarquía de títulos}
  \begin{tabularx}{\textwidth}{l l X}
    \toprule
    \headerrow
    \thead{Alias} & \thead{Color Ctp} & \thead{Razón} \\
    \midrule
    \texttt{doc-chapter}       & Pink      & Máximo contraste; warm; ancla el capítulo \\
    \texttt{doc-section}       & Sky       & Cool; segundo nivel; contrasta con Pink \\
    \texttt{doc-subsection}    & Lavender  & Suave; tercer nivel \\
    \texttt{doc-subsubsection} & Peach     & Warm suave; cuarto nivel \\
    \texttt{doc-paragraph}     & Subtext0  & Casi texto; diferencia mínima \\
    \bottomrule
  \end{tabularx}
\end{acadtable}

\subsection{Hipervínculos (\texttt{link-})}

Mapeo directo del Catppuccin Style Guide: Links=Blue, Followed=Lavender, Hover=Sky.

\begin{acadtable}[H]{ap:tab:link}{Aliases de hipervínculos}
  \begin{tabularx}{\textwidth}{l l X}
    \toprule
    \headerrow
    \thead{Alias} & \thead{Color Ctp} & \thead{Uso en hyperref} \\
    \midrule
    \texttt{link-internal} & Blue     & \texttt{linkcolor}: refs internas, TOC \\
    \texttt{link-external} & Peach    & \texttt{urlcolor}: URLs externas \\
    \texttt{link-cite}     & Green    & \texttt{citecolor}: citas bibliográficas \\
    \texttt{link-visited}  & Lavender & Followed Links \\
    \texttt{link-hover}    & Sky      & On Hover Links \\
    \bottomrule
  \end{tabularx}
\end{acadtable}

\subsection{Entornos teoremáticos (\texttt{env-})}

Cada entorno tiene \texttt{env-X} (frame) y \texttt{env-X-back} (fondo,
generado por \texttt{\textbackslash createback\{env-X\}}).

\begin{acadtable}[H]{ap:tab:env}{Aliases de entornos teoremáticos}
  \begin{tabularx}{\textwidth}{l l X}
    \toprule
    \headerrow
    \thead{Alias} & \thead{Color Ctp} & \thead{Entorno y razón} \\
    \midrule
    \texttt{env-def}  & Blue     & Definición: epistémico, fundacional \\
    \texttt{env-thm}  & Mauve    & Teorema: el más prominente del sistema \\
    \texttt{env-prop} & Green    & Proposición: verdad demostrada \\
    \texttt{env-ex}   & Peach    & Ejemplo: warm, invita a leer \\
    \texttt{env-obs}  & Overlay0 & Observación: secundaria, no compite \\
    \texttt{env-phys} & Sapphire & Conexión física: evoca ciencia \\
    \texttt{env-star} & Yellow   & Estrella Polar: máxima prominencia \\
    \texttt{env-comp} & Pink     & Ejercicio Computacional: activo \\
    \bottomrule
  \end{tabularx}
\end{acadtable}

\subsection{Matemáticas semánticas (\texttt{eq-})}

Tres ejes ortogonales. Ver \cref{ch:ecuaciones} para la documentación completa.

\begin{acadtable}[H]{ap:tab:eq}{Aliases de coloreado semántico de ecuaciones}
  \begin{tabularx}{\textwidth}{l l l X}
    \toprule
    \headerrow
    \thead{Alias} & \thead{Macro} & \thead{Color Ctp} & \thead{Eje y rol} \\
    \midrule
    \texttt{eq-op}      & \texttt{\textbackslash eqop}  & Blue      & Sintáctico: operadores \\
    \texttt{eq-fn}      & \texttt{\textbackslash eqfn}  & Green     & Sintáctico: funciones/mapas \\
    \texttt{eq-dm}      & \texttt{\textbackslash eqdm}  & Teal      & Sintáctico: diferenciales \\
    \texttt{eq-dom}     & \texttt{\textbackslash eqdom} & Sapphire  & Sintáctico: dominios/conjuntos \\
    \midrule
    \texttt{eq-unknown} & \texttt{\textbackslash equ}   & Peach     & Epistémico: incógnita \\
    \texttt{eq-known}   & \texttt{\textbackslash eqk}   & Mauve     & Epistémico: dato/fijado \\
    \texttt{eq-const}   & \texttt{\textbackslash eqcon} & Lavender  & Epistémico: constante universal \\
    \midrule
    \texttt{eq-vec}     & \texttt{\textbackslash eqvec} & Red       & Tipo: vector/campo \\
    \texttt{eq-ten}     & \texttt{\textbackslash eqten} & Rosewater & Tipo: tensor/matriz \\
    \bottomrule
  \end{tabularx}
\end{acadtable}

\subsection{Código fuente (\texttt{code-})}

Mapeo directo del Catppuccin Style Guide: Language Defaults (Code Editors).

\begin{acadtable}[H]{ap:tab:code}{Aliases de syntax highlighting}
  \begin{tabularx}{\textwidth}{l l X}
    \toprule
    \headerrow
    \thead{Alias} & \thead{Color Ctp} & \thead{Clase sintáctica} \\
    \midrule
    \texttt{code-keyword} & Mauve    & Keywords del lenguaje \\
    \texttt{code-string}  & Green    & Strings y literales \\
    \texttt{code-comment} & Overlay2 & Comentarios \\
    \texttt{code-operator}& Sky      & Operadores \\
    \texttt{code-type}    & Yellow   & Tipos / clases / interfaces \\
    \texttt{code-fn}      & Blue     & Funciones / métodos \\
    \texttt{code-mod}     & Teal     & Módulos / namespaces / enums \\
    \texttt{code-macro}   & Rosewater& Macros / decoradores / pragmas \\
    \texttt{code-builtin} & Red      & Builtins / atoms / symbols \\
    \texttt{code-const}   & Peach    & Constantes / números \\
    \texttt{code-param}   & Maroon   & Parámetros \\
    \texttt{code-delim}   & Overlay2 & Delimitadores (\texttt{\{\}[]()}) \\
    \texttt{code-bg}      & Base     & Fondo del bloque \\
    \texttt{code-border}  & Surface0 & Frame / regla \\
    \texttt{code-linenum} & Overlay1 & Números de línea \\
    \bottomrule
  \end{tabularx}
\end{acadtable}

\subsection{Algoritmos, plots, TikZ, tablas, captions y feedback}

Las tablas de aliases para los dominios \texttt{alg-}, \texttt{plot-},
\texttt{tikz-}, \texttt{table-}, \texttt{caption-} y \texttt{fb-}
se presentan en las secciones correspondientes de los capítulos:
\cref{ch:codigo} (algoritmos), \cref{ch:figuras} (plots y captions),
\cref{ch:diagramas} (TikZ), \cref{ch:tablas} (tablas).

%───────────────────────────────────────────────────────────────────────────────
\section{Protocolo de extensión}
\label{ap:sec:extension}
%───────────────────────────────────────────────────────────────────────────────

\begin{definicion}{Protocolo para añadir un dominio nuevo}
\label{ap:def:extension}
Cuando el template necesita soporte para un paquete o tipo de elemento
nuevo, el protocolo es exactamente este — no más, no menos:
\begin{enumerate}[label=\textbf{\arabic*.}, leftmargin=2em]
  \item Añadir una sección comentada en \texttt{themes/catppuccin-palette.tex}
    con los \texttt{\textbackslash colorlet} del nuevo dominio.
  \item Si los entornos necesitan fondos, llamar
    \texttt{\textbackslash createback\{alias\}} para cada uno.
  \item Crear \texttt{config/\{dominio\}.tex} consumiendo únicamente
    los aliases recién definidos. Nunca \texttt{Ctp*}.
  \item Registrar \texttt{\textbackslash input\{config/\{dominio\}.tex\}}
    en \texttt{main.tex} en el paso 8 (configs de dominio).
\end{enumerate}
\end{definicion}

El ejemplo canónico para un dominio \texttt{circuitikz} hipotético:

\begin{lstlisting}[style=latex, frame=none, numbers=none]
% En themes/catppuccin-palette.tex:
% - Circuitos (circuitikz) 
% Roles: pasivos, activos, fuentes, cables, labels
\colorlet{circ-passive}{CtpBlue}    % resistencias, capacitores
\colorlet{circ-active}{CtpGreen}    % transistores, diodos
\colorlet{circ-source}{CtpYellow}   % fuentes de voltaje/corriente
\colorlet{circ-wire}{CtpOverlay1}   % cables y conexiones
\colorlet{circ-label}{CtpSubtext0}  % labels de componentes

% En config/circuitikz.tex:
% Nunca Ctp* aqui
\ctikzset{
  resistors/scale=0.8,
  color=circ-wire,
}
\tikzset{
  component/.style={color=circ-passive, fill=circ-passive!10},
  source/.style   ={color=circ-source},
}

% En main.tex (paso 8):
\input{config/circuitikz.tex}
\end{lstlisting}

%───────────────────────────────────────────────────────────────────────────────
\section{El mecanismo \textbackslash createback}
\label{ap:sec:createback}
%───────────────────────────────────────────────────────────────────────────────

\begin{definicion}{\textbackslash createback}
\label{ap:def:createback}
El macro \texttt{\textbackslash createback\{alias\}} genera
\texttt{\{alias\}-back} como variante de fondo clara para
\texttt{tcolorbox}. La luminosidad se adapta automáticamente al flavor:
\begin{itemize}[leftmargin=2em]
  \item Tema claro (Latte):
    \texttt{\{alias\}-back} = \texttt{white!97!\{alias\}}
  \item Tema oscuro (Frappé, Macchiato, Mocha):
    \texttt{\{alias\}-back} = \texttt{CtpBase!95!\{alias\}}
\end{itemize}
\end{definicion}

Esto garantiza que todos los entornos teoremáticos tengan fondos
consistentes sin importar el flavor. El toggle \texttt{\textbackslash ifdarktheme}
se establece manualmente en \texttt{catppuccin-palette.tex} junto a la
línea de \texttt{\textbackslash usepackage[Flavor]}.

%───────────────────────────────────────────────────────────────────────────────
\section{Fundamento en el Style Guide oficial de Catppuccin}
\label{ap:sec:fundamento}
%───────────────────────────────────────────────────────────────────────────────

Las asignaciones de este template no son arbitrarias. Cada color de dominio
tiene una correspondencia directa con las categorías del Style Guide oficial
de Catppuccin. La siguiente tabla traza esa correspondencia para los dominios
de mayor visibilidad:

\begin{acadtable}[H]{ap:tab:sg-mapping}{Correspondencia con el Style Guide oficial}
  \begin{tabularx}{\textwidth}{l l X}
    \toprule
    \headerrow
    \thead{Alias del template} & \thead{Categoría Style Guide} & \thead{Descripción SG} \\
    \midrule
    \texttt{link-internal}   & Links, URLs      & ``Blue for links'' \\
    \texttt{link-visited}    & Followed Links   & ``Lavender for followed links'' \\
    \texttt{link-hover}      & On Hover Links   & ``Sky for hover'' \\
    \texttt{code-keyword}    & Keyword          & ``Mauve for keywords'' \\
    \texttt{code-string}     & Strings          & ``Green for strings'' \\
    \texttt{code-comment}    & Comments         & ``Overlay 2 for comments'' \\
    \texttt{code-operator}   & Operators        & ``Sky for operators'' \\
    \texttt{code-type}       & Classes, Types   & ``Yellow for types'' \\
    \texttt{code-fn}         & Methods          & ``Blue for methods'' \\
    \texttt{fb-success}      & Success          & ``Green for success'' \\
    \texttt{fb-warning}      & Warnings         & ``Yellow for warnings'' \\
    \texttt{fb-error}        & Errors           & ``Red for errors'' \\
    \texttt{fb-info}         & Information      & ``Teal for information'' \\
    \texttt{env-thm}         & Highest priority & ``Mauve for most prominent'' \\
    \texttt{rainbow-1..6}    & Rainbow Highlights & Secuencia de colores ordenada \\
    \bottomrule
  \end{tabularx}
\end{acadtable}

Las extensiones propias del template (dominios \texttt{eq-}, \texttt{doc-},
\texttt{tikz-}, \texttt{table-}, \texttt{caption-}, \texttt{plot-}) no
tienen correspondencia directa en el Style Guide porque ese no cubre
documentos matemáticos. Las asignaciones se derivaron de los principios
del Style Guide: colores fríos para elementos epistémicos neutros (Blue
para operadores, Sapphire para dominios), colores cálidos para elementos
con carga semántica fuerte (Peach para incógnitas, Red para vectores).

\begin{notasbib}
El Style Guide oficial de Catppuccin está disponible en
\texttt{https://github.com/catppuccin/catppuccin/blob/main/docs/style-guide.md}.
La paleta completa con valores hexadecimales para cada flavor se documenta
en \texttt{https://catppuccin.com/palette}. El paquete
\texttt{catppuccinpalette} para \LaTeX\ se mantiene en CTAN y en
\texttt{https://github.com/catppuccin/latex}.
\end{notasbib}

  %═══════════════════════════════════════════════════════════════════════════════
\chapter{Galería de código fuente}
\label{ap:codigo-galeria}
%═══════════════════════════════════════════════════════════════════════════════

Este apéndice presenta ejemplos de código en los diez lenguajes configurados
en la plantilla. Cada estilo mapea las categorías sintácticas del lenguaje a
la paleta Catppuccin según la tabla del \cref{sec:codigo-breve}. Los estilos
se invocan con \verb|style=python|, \verb|style=rust|, etc.

\dificultad{1}{1}

%───────────────────────────────────────────────────────────────────────────────
\section{Python}
%───────────────────────────────────────────────────────────────────────────────

\begin{lstlisting}[style=python, caption={Descenso de gradiente con momentum de Nesterov}]
import numpy as np
from typing import Callable, Tuple

class SGDMomentum:
    """SGD optimizer with Nesterov momentum."""

    def __init__(self, lr: float = 1e-3, momentum: float = 0.9):
        self.lr = lr
        self.momentum = momentum
        self.velocity = None

    def step(self, params: np.ndarray, grad: np.ndarray) -> np.ndarray:
        if self.velocity is None:
            self.velocity = np.zeros_like(params)
        # Nesterov momentum update
        self.velocity = self.momentum * self.velocity - self.lr * grad
        return params + self.momentum * self.velocity - self.lr * grad

def train(model: Callable, X: np.ndarray, y: np.ndarray,
          epochs: int = 100) -> list:
    optimizer = SGDMomentum(lr=0.01)
    losses = []
    for epoch in range(epochs):
        y_hat = model(X)
        loss = np.mean((y_hat - y) ** 2)
        losses.append(loss)
        if loss < 1e-6:
            print(f"Converged at epoch {epoch}")
            break
    return losses
\end{lstlisting}

%───────────────────────────────────────────────────────────────────────────────
\section{C}
%───────────────────────────────────────────────────────────────────────────────

\begin{lstlisting}[style=c, caption={Producto matriz-vector}]
#include <stdio.h>
#include <stdlib.h>
#include <string.h>
#include <math.h>

void matvec(const double *A, const double *x, double *y,
            size_t m, size_t n) {
    memset(y, 0, m * sizeof(double));
    for (size_t i = 0; i < m; i++)
        for (size_t j = 0; j < n; j++)
            y[i] += A[i * n + j] * x[j];
}

double norm2(const double *v, size_t n) {
    double s = 0.0;
    for (size_t i = 0; i < n; i++)
        s += v[i] * v[i];
    return sqrt(s);
}

int main(void) {
    double A[] = {1.0, 2.0, 3.0, 4.0, 5.0, 6.0};
    double x[] = {1.0, 1.0, 1.0};
    double y[2];
    matvec(A, x, y, 2, 3);
    printf("||y|| = %.6f\n", norm2(y, 2));
    return EXIT_SUCCESS;
}
\end{lstlisting}

%───────────────────────────────────────────────────────────────────────────────
\section{C++}
%───────────────────────────────────────────────────────────────────────────────

\begin{lstlisting}[style=cpp, caption={Integrador de Monte Carlo con templates}]
#include <iostream>
#include <vector>
#include <random>
#include <functional>

template<typename T>
class MonteCarloIntegrator {
private:
    std::mt19937 rng_;
    std::uniform_real_distribution<T> dist_;
public:
    explicit MonteCarloIntegrator(unsigned seed = 42)
        : rng_(seed), dist_(0.0, 1.0) {}

    T integrate(std::function<T(T)> f, T a, T b, size_t N) {
        T sum = 0.0;
        for (size_t i = 0; i < N; ++i) {
            T x = a + (b - a) * dist_(rng_);
            sum += f(x);
        }
        return (b - a) * sum / static_cast<T>(N);
    }
};

int main() {
    MonteCarloIntegrator<double> mc;
    auto result = mc.integrate(
        [](double x) { return std::sin(x); },
        0.0, M_PI, 1'000'000
    );
    std::cout << "Result: " << result << std::endl;
    return 0;
}
\end{lstlisting}

%───────────────────────────────────────────────────────────────────────────────
\section{Rust}
%───────────────────────────────────────────────────────────────────────────────

\begin{lstlisting}[style=rust, caption={Solucionador RK4 con traits}]
trait OdeSolver {
    fn step(&self, f: &dyn Fn(f64, &[f64]) -> Vec<f64>,
            t: f64, y: &[f64], dt: f64) -> Vec<f64>;
}

struct RungeKutta4;

impl OdeSolver for RungeKutta4 {
    fn step(&self, f: &dyn Fn(f64, &[f64]) -> Vec<f64>,
            t: f64, y: &[f64], dt: f64) -> Vec<f64> {
        let n = y.len();
        let k1 = f(t, y);
        let y2: Vec<f64> = (0..n).map(|i| y[i] + 0.5*dt*k1[i]).collect();
        let k2 = f(t + 0.5*dt, &y2);
        let y3: Vec<f64> = (0..n).map(|i| y[i] + 0.5*dt*k2[i]).collect();
        let k3 = f(t + 0.5*dt, &y3);
        let y4: Vec<f64> = (0..n).map(|i| y[i] + dt*k3[i]).collect();
        let k4 = f(t + dt, &y4);
        (0..n).map(|i| {
            y[i] + dt/6.0 * (k1[i] + 2.0*k2[i] + 2.0*k3[i] + k4[i])
        }).collect()
    }
}

fn main() {
    let rk4 = RungeKutta4;
    let harmonic = |_t: f64, y: &[f64]| -> Vec<f64> {
        vec![y[1], -y[0]]
    };
    let mut y = vec![1.0, 0.0];
    for i in 0..1000 {
        y = rk4.step(&harmonic, i as f64 * 0.01, &y, 0.01);
    }
    println!("y(10) = {:.6}, y'(10) = {:.6}", y[0], y[1]);
}
\end{lstlisting}

%───────────────────────────────────────────────────────────────────────────────
\section{Go}
%───────────────────────────────────────────────────────────────────────────────

\begin{lstlisting}[style=go, caption={Problema de Basilea con goroutines}]
package main

import (
    "fmt"
    "math"
    "sync"
)

type Result struct {
    Index int
    Value float64
}

func partialSum(start, end int, ch chan<- Result, wg *sync.WaitGroup) {
    defer wg.Done()
    sum := 0.0
    for k := start; k < end; k++ {
        sum += 1.0 / math.Pow(float64(k), 2)
    }
    ch <- Result{Index: start, Value: sum}
}

func main() {
    const N = 1_000_000
    const workers = 8
    chunk := N / workers

    ch := make(chan Result, workers)
    var wg sync.WaitGroup

    for i := 0; i < workers; i++ {
        wg.Add(1)
        s := i*chunk + 1
        e := s + chunk
        if i == workers-1 { e = N + 1 }
        go partialSum(s, e, ch, &wg)
    }
    go func() { wg.Wait(); close(ch) }()

    total := 0.0
    for r := range ch { total += r.Value }
    fmt.Printf("pi ~ %.10f\n", math.Sqrt(6.0*total))
}
\end{lstlisting}

%───────────────────────────────────────────────────────────────────────────────
\section{Julia}
%───────────────────────────────────────────────────────────────────────────────

\begin{lstlisting}[style=julia, caption={Euler--Maruyama con dispatch múltiple}]
using Random, LinearAlgebra, Statistics

abstract type AbstractSDE end

struct OrnsteinUhlenbeck <: AbstractSDE
    theta::Float64
    mu::Float64
    sigma::Float64
end

drift(sde::OrnsteinUhlenbeck, x, t) = sde.theta * (sde.mu - x)
diffusion(sde::OrnsteinUhlenbeck, x, t) = sde.sigma

function euler_maruyama(sde::AbstractSDE, x0::Float64,
                        dt::Float64, T::Float64; seed::Int=42)
    rng = MersenneTwister(seed)
    N = Int(T / dt)
    x = zeros(N + 1)
    x[1] = x0
    @inbounds for i in 1:N
        dW = sqrt(dt) * randn(rng)
        x[i+1] = x[i] + drift(sde, x[i], (i-1)*dt)*dt +
                         diffusion(sde, x[i], (i-1)*dt)*dW
    end
    return range(0.0, T; length=N+1), x
end

ou = OrnsteinUhlenbeck(2.0, 0.0, 1.0)
t, x = euler_maruyama(ou, 1.0, 1e-3, 10.0)
@info "Stationary variance" exact=ou.sigma^2/(2*ou.theta) numerical=var(x[end-1000:end])
\end{lstlisting}

%───────────────────────────────────────────────────────────────────────────────
\section{MATLAB}
%───────────────────────────────────────────────────────────────────────────────

\begin{lstlisting}[style=matlab, caption={Elementos finitos 1D}]
function [u, x] = fem1d(N, f_rhs, a, b)
% FEM1D  Solves -u''(x) = f(x) on [a,b], u(a)=u(b)=0
%        Linear finite elements with N elements.

h = (b - a) / N;
x = linspace(a, b, N+1)';
K = sparse(N+1, N+1);
F = zeros(N+1, 1);

for e = 1:N
    Ke = (1/h) * [1 -1; -1 1];
    xm = (x(e) + x(e+1)) / 2;
    Fe = h/2 * [f_rhs(xm); f_rhs(xm)];
    idx = [e, e+1];
    K(idx, idx) = K(idx, idx) + Ke;
    F(idx) = F(idx) + Fe;
end

interior = 2:N;
u = zeros(N+1, 1);
u(interior) = K(interior, interior) \ F(interior);

u_exact = @(x) sin(pi*x) / pi^2;
err = sqrt(h * sum((u - u_exact(x)).^2));
fprintf('L2 error: %.2e (h = %.2e)\n', err, h);
end
\end{lstlisting}

%───────────────────────────────────────────────────────────────────────────────
\section{Fortran}
%───────────────────────────────────────────────────────────────────────────────

\begin{lstlisting}[style=fortran, caption={Factorización LU con pivoteo parcial}]
module linear_algebra
  implicit none
  integer, parameter :: dp = selected_real_kind(15, 307)
contains
  subroutine lu_factor(A, n, ipiv, info)
    integer, intent(in)     :: n
    real(dp), intent(inout) :: A(n, n)
    integer, intent(out)    :: ipiv(n), info
    integer  :: i, j, k, imax
    real(dp) :: maxval, temp

    info = 0
    do k = 1, n
      maxval = abs(A(k, k));  imax = k
      do i = k + 1, n
        if (abs(A(i, k)) > maxval) then
          maxval = abs(A(i, k));  imax = i
        end if
      end do
      ipiv(k) = imax
      if (maxval < epsilon(1.0_dp)) then
        info = k;  return
      end if
      if (imax /= k) then
        do j = 1, n
          temp = A(k,j); A(k,j) = A(imax,j); A(imax,j) = temp
        end do
      end if
      do i = k + 1, n
        A(i,k) = A(i,k) / A(k,k)
        do j = k + 1, n
          A(i,j) = A(i,j) - A(i,k) * A(k,j)
        end do
      end do
    end do
  end subroutine lu_factor
end module linear_algebra
\end{lstlisting}

%───────────────────────────────────────────────────────────────────────────────
\section{Assembler (x86-64)}
%───────────────────────────────────────────────────────────────────────────────

\begin{lstlisting}[style=asm, caption={Producto punto con SSE}]
; dot_product_sse(float *a, float *b, int n) -> float
; System V AMD64: rdi=a, rsi=b, edx=n; returns in xmm0
section .text
global dot_product_sse

dot_product_sse:
    xorps   xmm0, xmm0
    test    edx, edx
    jle     .done
    mov     ecx, edx
    shr     ecx, 2
    jz      .remainder
.loop4:
    movups  xmm1, [rdi]
    movups  xmm2, [rsi]
    mulps   xmm1, xmm2
    addps   xmm0, xmm1
    add     rdi, 16
    add     rsi, 16
    dec     ecx
    jnz     .loop4
    ; horizontal reduction
    movaps  xmm1, xmm0
    shufps  xmm1, xmm1, 0x4E
    addps   xmm0, xmm1
    movaps  xmm1, xmm0
    shufps  xmm1, xmm1, 0xB1
    addss   xmm0, xmm1
.remainder:
    and     edx, 3
    jz      .done
.loop1:
    movss   xmm1, [rdi]
    mulss   xmm1, [rsi]
    addss   xmm0, xmm1
    add     rdi, 4
    add     rsi, 4
    dec     edx
    jnz     .loop1
.done:
    ret
\end{lstlisting}

%───────────────────────────────────────────────────────────────────────────────
\section{Mathematica}
%───────────────────────────────────────────────────────────────────────────────

\begin{lstlisting}[style=mathematica, caption={EDP de calor: simbólica + numérica}]
(* 1D heat equation: du/dt = k * d^2u/dx^2 *)

k = 0.01;  L = Pi;

uExact[x_, t_, nTerms_] := Module[{n},
    Sum[(2/(n*Pi)) * Sin[n*x] * Exp[-k*n^2*t],
        {n, 1, nTerms, 2}]
]

sol = NDSolve[{
    D[u[x,t], t] == k * D[u[x,t], {x,2}],
    u[0,t] == 0,  u[L,t] == 0,
    u[x,0] == If[x < L/2, 2*x/L, 2*(1 - x/L)]
}, u, {x, 0, L}, {t, 0, 10}];

error = NIntegrate[
    (u[x,5] /. sol[[1]] - uExact[x, 5, 50])^2,
    {x, 0, L}] // Sqrt;

Print["L2 error at t=5: ", ScientificForm[error, 3]]
\end{lstlisting}

%───────────────────────────────────────────────────────────────────────────────
\section*{Referencia rápida de estilos}
\addcontentsline{toc}{section}{Referencia rápida de estilos}
%───────────────────────────────────────────────────────────────────────────────

\begin{acadtable}{tab:estilos-listings}{Estilos de listings disponibles}
  \begin{tabularx}{\textwidth}{l l Y}
    \toprule
    \headerrow
    \thead{Estilo} & \thead{Lenguaje} & \thead{Uso} \\
    \midrule
    \texttt{python}      & Python         & \verb|style=python| \\
    \rowcolor{table-row-alt}
    \texttt{c}           & C              & \verb|style=c| \\
    \texttt{cpp}         & C++            & \verb|style=cpp| \\
    \rowcolor{table-row-alt}
    \texttt{rust}        & Rust           & \verb|style=rust| \\
    \texttt{go}          & Go             & \verb|style=go| \\
    \rowcolor{table-row-alt}
    \texttt{julia}       & Julia          & \verb|style=julia| \\
    \texttt{matlab}      & MATLAB/Octave  & \verb|style=matlab| \\
    \rowcolor{table-row-alt}
    \texttt{fortran}     & Fortran (2008) & \verb|style=fortran| \\
    \texttt{asm}         & x86-64 (Intel) & \verb|style=asm| \\
    \rowcolor{table-row-alt}
    \texttt{mathematica} & Wolfram        & \verb|style=mathematica| \\
    \bottomrule
  \end{tabularx}
\end{acadtable}

\begin{notasbib}
La configuración de \texttt{listings} sigue la guía de estilo oficial de
Catppuccin (\url{https://catppuccin.com/style-guide}). Los lenguajes Rust
y Go se definen desde cero con \verb|\lstdefinelanguage| porque el paquete
\texttt{listings} no los incluye nativamente.
Para resaltado más avanzado (funciones definidas por el usuario, type
inference), considerar \texttt{minted} con Pygments como alternativa.
\end{notasbib}

\begin{notasbib}
La configuración de \texttt{listings} sigue la guía de estilo oficial de
Catppuccin. Los lenguajes Rust y Go se definen desde cero con
\texttt{\textbackslash lstdefinelanguage} porque \texttt{listings}
no los incluye nativamente.
\end{notasbib}

\end{appendices}

% En appendixA.tex:
\chapter{Style Guide de Catppuccin Académico}
\label{ap:styleguide}
\end{lstlisting}

Las referencias a apéndices usan el mismo sistema que cualquier otro
elemento: \texttt{\textbackslash cref\{ap:styleguide\}} produce
\textit{ap.~A} automáticamente, con el nombre registrado en
\texttt{config/packages.tex}.

\begin{observacion}{Numeración de figuras y tablas en apéndices}
Por defecto, las figuras en el apéndice A se numeran A.1, A.2, etc.
Este comportamiento es automático con el paquete \texttt{appendix}.
No se necesita configuración adicional.
\end{observacion}

%───────────────────────────────────────────────────────────────────────────────
\section{Encabezados y pies de página}
\label{sec:headers}
%───────────────────────────────────────────────────────────────────────────────

El template usa \texttt{fancyhdr}~\cite{fancyhdr2022} configurado en
\texttt{config/geometry.tex}. El estilo por defecto muestra el nombre de
la sección actual a la izquierda y el número de página a la derecha.
Las páginas de inicio de capítulo usan el estilo \texttt{plain}
(solo número de página centrado, sin línea de encabezado).

Para modificar el contenido de los encabezados, editar directamente
\texttt{config/geometry.tex}:

\begin{lstlisting}[style=latex, frame=none, numbers=none]
\fancyhead[L]{\nouppercase{\leftmark}}   % seccion actual, izquierda
\fancyhead[R]{\thepage}                  % numero de pagina, derecha
\fancyhead[C]{}                          % centro vacio
\fancyfoot[C]{}                          % pie vacio
\end{lstlisting}

En documentos de doble cara (\texttt{twoside}), las posiciones L y R se
alternan entre páginas pares e impares. Para ese caso, usar las posiciones
\texttt{LE}, \texttt{RO} (par izquierda, impar derecha) en lugar de
\texttt{L} y \texttt{R}.

%───────────────────────────────────────────────────────────────────────────────
\printbibliography[heading=subbibliography]
